\chapter{Problem Statement: Mixed Cache Traffic on One Link}
\label{ch:problem}

When inference is split between a prefill server and a decode server, a
single link between them carries everything: hundreds of megabytes of cache,
and messages of a few bytes that tell the other side that the cache is ready.
An inventory of that traffic comes first. Arithmetic and a queueing model
then show why the small messages suffer when they share a queue with the
large ones, existing systems are examined for how they treat the two kinds,
and the research questions follow from what those systems leave open.

\section{Traffic on the Link}
\label{sec:traffic}

\Cref{ch:inference} showed that the cache built during prefill is a pure
function of the prompt and the weights, so it can be built on one machine and
used on another. In Granite 3.3 2B \citep{granite33} each layer stores a
key vector and a value vector per token for each of its key-value heads, so
the cache of one layer and the cache of the whole model occupy
\mathalt{s layer equals 2 times h times d times w times T, and s cache equals L times
  s layer}{%
\begin{equation}
  s_{\text{layer}} = 2\, h\, d\, w\, T
  \qquad\text{and}\qquad
  s_{\text{cache}} = L\, s_{\text{layer}}
  \label{eq:kvsize}
\end{equation}}
bytes, where the factor $2$ counts the key and the value, $h = 8$ is the
number of key-value heads of a layer, $d = 64$ is the number of entries of
each key and each value, $w = 2$ is the number of bytes of one entry in the
16-bit format, $T$ is the number of tokens of the prompt, and $L = 40$ is the
number of layers. One token therefore occupies
$2 \times 8 \times 64 \times 2 = 2{,}048$ bytes, or 2 KiB, in each layer and
80 KiB over all layers, and a prompt of $T = 2{,}048$ tokens produces 4 MiB
per layer and 160 MiB in total. On the laboratory link the payload alone
needs the time
\mathalt{t wire equals 8 times s cache divided by C, which is 8 times 160 times 2 to the 20
  bits, divided by 10 to the 10 bits per second, which equals 134.2 milliseconds}{%
\begin{equation}
  t_{\text{wire}} = \frac{8\, s_{\text{cache}}}{C}
  = \frac{8 \times 160 \times 2^{20}\ \text{bits}}{10^{10}\ \text{bits/s}} = 134.2\ \text{ms},
  \label{eq:wiretime}
\end{equation}}
where $C$ is the rate of the link in bits per second, here 10 gigabits per
second, and the factor $8$ converts bytes to bits.
Each RoCEv2 packet adds 82 bytes of headers, trailer and inter-frame gap, so
the wire time is about 137 ms with 4,096-byte RDMA packets and 145 ms with
1,024-byte packets. These are lower bounds from arithmetic, not measurements.

The cache is not the only traffic on the link. \Cref{tab:traffic} lists the
classes that occur in split serving, and \cref{sec:overview} states which of
them the connector built here sends.

\begin{table}[htbp]
  \centering
  \caption[Traffic between a prefill server and a decode server]{Traffic between a prefill server and a decode server in split serving, with sizes for Granite 3.3 2B.}
  \label{tab:traffic}
  \small
  \begin{tabular}{L{4.1cm}L{3.6cm}L{6.4cm}}
    \toprule
    Message & Size & What waits on it \\
    \midrule
    Prefill cache, one layer & 4 MiB for 2,048 tokens & nothing until the whole request has arrived \\
    End-of-request marker & under 100 bytes & the decode server, in a design that starts on its arrival \\
    Flow-control credit & 8 bytes & the sender's next block \\
    Token identifiers and request metadata & tens of bytes & the request itself \\
    Decode-side cache update & 2 KiB per token per layer & the next token of a running request \\
    \bottomrule
  \end{tabular}
\end{table}

The classes differ in two ways that matter. The large ones are measured in
megabytes and no part of the system waits on any single block. The small ones are
measured in bytes and some part of the system is blocked until each arrives. We call the first kind cold and the second kind hot.

\section{Head-of-Line Blocking}
\label{sec:hol}

If hot and cold messages share one first-in first-out queue, a hot message
that arrives behind $k$ cold messages waits for all of them, for the time
\mathalt{W equals 8 over C, times the sum of s i for i from 1 to k}{%
\begin{equation}
  W = \frac{8}{C} \sum_{i=1}^{k} s_i ,
  \label{eq:fifo}
\end{equation}}
where $s_i$ is the size of the $i$th cold message in bytes and $C$ is the
rate of the link as before, so that each term $8 s_i / C$ is the time for
which one of these messages occupies the link.
Behind four 1 MiB blocks at that rate the wait is 3.36 ms, for a
message whose own transmission takes a few microseconds. \Cref{fig:hol} shows
the two arrangements.

\begin{figure}[htbp]
  \centering
  \begin{tikzpicture}[diagram,
      alt={Two panels compare two ways to queue messages for one link. Panel a, one shared
        queue: a large cold block is at the head of the queue and three small hot messages wait
        behind it. Panel b, one queue per class: the three hot messages sit in their own queue
        and the cold block sits in a second queue, so the hot messages reach the link without
        waiting for the block. Hot messages are small plain squares; the cold block is a wide
        hatched rectangle.},
      panel/.style={draw=dGray, fill=dGray!6, line width=0.6pt, rounded corners=6pt},
      queue/.style={draw=dInk, fill=white, line width=0.8pt, rounded corners=3pt},
      small item/.style={draw=dBlue, fill=dBlue!12, line width=0.8pt},
      large item/.style={draw=dOrange, fill=white, line width=0.8pt,
                         postaction={pattern=north east lines, pattern color=dOrange}},
      link/.style={tint=dGreen, minimum width=13mm, minimum height=13mm, inner sep=2pt}]
    \draw[panel] (-0.4,-2.0) rectangle (6.95,2.3);
    \node[heading, anchor=west] at (-0.1,1.85) {(a) One shared queue};
    \draw[queue] (0,-0.35) rectangle (4.7,0.65);
    \foreach \x in {0.35,1.05,1.75} { \draw[small item] (\x,-0.13) rectangle ++(0.56,0.56); }
    \draw[large item] (2.5,-0.13) rectangle (4.45,0.43);
    \node[link] (linkA) at (6.0,0.15) {Link};
    \draw[flow] (4.7,0.15) -- (linkA.west);
    \node[small text, anchor=west] at (-0.1,-1.55) {Hot messages wait behind the block.};
    \begin{scope}[xshift=7.85cm]
      \draw[panel] (-0.4,-2.0) rectangle (6.95,2.3);
      \node[heading, anchor=west] at (-0.1,1.85) {(b) One queue per class};
      \draw[queue] (0,0.35) rectangle (4.7,1.35);
      \foreach \x in {2.49,3.19,3.89} { \draw[small item] (\x,0.57) rectangle ++(0.56,0.56); }
      \draw[queue] (0,-1.05) rectangle (4.7,-0.05);
      \draw[large item] (2.5,-0.83) rectangle (4.45,-0.27);
      \node[link] (linkB) at (6.0,0.15) {Link};
      \draw[flow] (4.7,0.85)  -- (linkB.north west);
      \draw[flow] (4.7,-0.55) -- (linkB.south west);
      \node[small text, anchor=west] at (-0.1,-1.55) {Hot messages do not wait for the block.};
    \end{scope}
    \draw[small item] (4.3,-2.85) rectangle ++(0.42,0.42);
    \node[small text, anchor=west] at (4.85,-2.64) {Hot message};
    \draw[large item] (7.6,-2.85) rectangle ++(1.1,0.42);
    \node[small text, anchor=west] at (8.85,-2.64) {Cold block};
  \end{tikzpicture}
  \caption[Head-of-line blocking in a shared queue and with one queue per class]{Head-of-line blocking in a shared queue, and the same traffic with one queue per class.}
  \label{fig:hol}
\end{figure}

A queueing model gives the average case. Treat the link as one server at
which messages arrive at random, as a Poisson process, and are sent one at a
time. The mean wait in a first-in first-out queue is then given by the
Pollaczek-Khinchine formula,
\mathalt{W bar equals lambda times the expected value of S squared, divided by 2 times
  1 minus rho}{%
\begin{equation}
  \overline{W} = \frac{\lambda\, \mathrm{E}[S^2]}{2(1-\rho)},
  \label{eq:pk}
\end{equation}}
where $\overline{W}$ is the mean wait of a message before its transmission
begins, $\lambda$ is the arrival rate in messages per second, $S$ is the
service time, which is $8s/C$ seconds for a message of $s$ bytes,
$\mathrm{E}[S^2]$ is the mean of its square, and
$\rho = \lambda\, \mathrm{E}[S]$ is the fraction of time for which the link is
busy. The second moment is the damaging term: a few very large messages raise
the wait of every message. With
two queues and a rule that the hot queue is served first whenever the link
becomes free, the hot wait drops to the residual of whatever is in service
\citep{cobham1954}; and if hot packets can be interleaved between the packets
of a cold message, it drops again.

\Cref{tab:holmodel} evaluates the three cases for one illustrative load. The load and the Poisson arrivals are assumptions, so the numbers are model
output: they show the size of the effect that \cref{sec:res-e3} is to measure
and are not that measurement.

\begin{table}[htbp]
  \centering
  \caption[Mean wait of a hot message in a queueing model]{Mean wait of a hot message in a queueing model. Assumed load: a 10 Gb/s link; cold 1 MiB messages at 60 percent load; hot 4 KiB messages at 10,000 per second.}
  \label{tab:holmodel}
  \small
  \begin{tabular}{L{8.6cm}L{3.2cm}}
    \toprule
    Arrangement & Mean hot wait (model) \\
    \midrule
    One first-in first-out queue for both classes & 0.69 ms \\
    Two queues, hot served first, cold messages not interrupted & 0.26 ms \\
    Two queues, hot packets interleaved between cold packets & 1.07\us \\
    \bottomrule
  \end{tabular}
\end{table}

The design presented here guarantees the step from the first row to the
second: with a separate queue pair per class, no cold message is ever ahead of
a hot one in its own send queue. The step from the second row to the third depends on what sits below the
queue pairs: a hardware network card arbitrates between queue pairs packet by
packet, whereas the software RoCE driver processes all queue pairs of a
device in kernel context, and whether it interleaves them finely enough is
the open question that experiment E3 answers.

\section{Existing Systems}

Disaggregated serving is established. DistServe \citep{zhong2024distserve} and
Splitwise \citep{patel2024splitwise} showed that prefill and decode have
different bottlenecks and gain from running on separate machines, and Mooncake
\citep{qin2024mooncake} built a serving architecture around the cache. The
open-source engine used here, vLLM \citep{kwon2023vllm,vllm_repo}, supports
the split through a connector interface, with NVIDIA's NIXL \citep{nixl} as
the usual transport, and LMCache \citep{lmcache_blog} and llm-d \citep{llmd}
build on the same engine. A review of these systems and their documentation identified five gaps.

First, the transports are built for bulk. They move a cache as fast as
possible, and none of them distinguishes the 8-byte credit from the 4 MiB
block when choosing a queue. Second, the portable path is shaped for one kind
of adapter. In the NIXL revision read for this work, the libfabric back end
looks first for the provider of HPE's Slingshot network, then for that of a
cloud vendor's adapter, and then falls back to TCP \citep{nixl}; the verbs
provider that would serve RoCE on commodity cards is not in its list, and its
documentation asks for a newer libfabric and for GPUDirect. Third, transport
libraries have side effects. An engineering report from a model provider
traced a memory leak of hundreds of megabytes per minute in a production
split-serving deployment to the memory hooks of the transport library under
the connector \citep{mistral2026leak}. Fourth, RDMA paths fail quietly. A
production parallel file system that uses RDMA has reported such issues,
including a race that returned incorrect file data \citep{ibm_ij44629} and a
reconnection fault \citep{ibm_iv61625}. Fifth, completions go unchecked. A
benchmark that counts completions without reading their status records a
failed operation as a very fast one (\cref{sec:rules}).

Work on the collective side of GPU communication has reached the same
conclusion from the other direction: for small messages the time is in the
software path and the queueing, not on the wire \citep{shen2026sol}. In addition, a library for point-to-point RDMA in language-model systems,
fabric-lib \citep{licker2026fabriclib}, treats the completion signal as a first-class design choice.

\section{The Gap and the Research Questions}
\label{sec:rq}

No system studied here treats the traffic between the two halves of a split
inference server as mixed, on commodity Ethernet with software RoCE and with the cache staged through host memory. That
setting is the subject of this work, and the intended contribution is the
classification and scheduling of mixed traffic, not a faster bulk path.
Three questions follow.

\begin{enumerate}
  \item[RQ1.] \rqone{} (E3, E4)
  \item[RQ2.] \rqtwo{} (E1, E2, E5)
  \item[RQ3.] \rqthree{} (E6, E7, E8)
\end{enumerate}

The first half of RQ3 is measured on the serving path, whereas the decisions
named in its second half are made by components that are evaluated
separately (\cref{sec:overview}). All three questions concern what happens
between a program and its network card, so the next chapter describes that
interface from the bottom up.
